% OS固有のフォントに依存しないTeX Live収録フォント。
\usepackage[haranoaji]{luatexja-preset}
\usepackage{fontspec}
\setmainfont{Latin Modern Roman}
\setsansfont{Latin Modern Sans}
\setmonofont{Latin Modern Mono}
\usepackage{amsmath,amssymb,amsthm}
\usepackage{graphicx,booktabs}
\usepackage{xcolor}
\usepackage{listings}
\usepackage{url}
\usepackage[unicode,hidelinks]{hyperref}
\hypersetup{pdftitle={\BookTitle},pdfauthor={\CircleName}}
\graphicspath{{images/}}

\setlength{\parindent}{1\zw}
\setlength{\parskip}{0pt}
\setcounter{secnumdepth}{2}
\setcounter{tocdepth}{2}
\ModifyHeading{section}{font={\large\sffamily\bfseries}}
\ModifyHeading{subsection}{font={\normalsize\sffamily\bfseries}}
\NewPageStyle{techbook}{
  running_head_position=top-fore-edge,
  odd_running_head=_chapter,even_running_head={\BookTitle},
  mark_format={_chapter={#1}},
  running_head_font={\small},
  nombre_position=bottom-fore-edge,nombre_font={\small}}
% 章頭もページ番号を表示する。章見出し自体はjlreq標準を使う。
\ModifyPageStyle{plain}{nombre_position=bottom-fore-edge}
\newcommand{\chapterauthor}[1]{%
  \begin{flushright}\small 執筆：#1\end{flushright}}

% 本文はモノクロ。表紙だけ色を使用できる。
\definecolor{CoverAccent}{RGB}{28,66,100}
\renewcommand{\lstlistingname}{コード}
\lstset{
  language=Python,basicstyle=\small\ttfamily,
  keywordstyle=\normalfont\small\sffamily\bfseries,commentstyle=\itshape,stringstyle=\ttfamily,
  numbers=left,numberstyle=\scriptsize,numbersep=8pt,
  frame=single,breaklines=true,breakatwhitespace=false,
  columns=fullflexible,keepspaces=true,showstringspaces=false,
  tabsize=4,captionpos=b,xleftmargin=2em,framexleftmargin=0.5em,
  aboveskip=\baselineskip,belowskip=\baselineskip}
% listingsのコード内はASCII推奨。日本語の解説は本文・キャプションに書く。
\theoremstyle{definition}
\newtheorem{definition}{定義}[chapter]
\theoremstyle{plain}
\newtheorem{theorem}{定理}[chapter]
\renewcommand{\proofname}{証明}
